% Part: computability
% Chapter: computability-theory
% Section: normal-form

\documentclass[../../../include/open-logic-section]{subfiles}

\begin{document}

\olfileid[pt]{cmp}{thy}{nfm}
\olsection{O teorema da forma normal}

Suponha que possamos descrever definições de funções computáveis e
testar se uma suposta descrição do registro completo da computação de
uma dessas funções em certa entrada está correta. Então é razoável
concluir que, independentemente do modelo de computação, podemos
determinar o valor de qualquer função computável~$f$ em qualquer
entrada~$x$ da seguinte maneira:
\begin{enumerate}
  \item Percorra todas as descrições possíveis de registros de
  computação.
  \item Teste se determinado registro é o registro de uma computação
  de~$f(x)$.
  \item Extraia do registro correto o valor de~$f(x)$, caso o tenhamos
  encontrado.
\end{enumerate}
O teorema da forma normal de Kleene afirma que isso de fato é verdade.

\begin{thm}[Teorema da forma normal de Kleene]
\ollabel{thm:normal-form}
Existe uma relação recursiva primitiva~$T(e, x, s)$ e uma função
recursiva primitiva~$U(s)$ com a seguinte propriedade: se $f$ é uma
função parcialmente computável qualquer, então, para algum~$e$,
\[
f(x) \simeq U(\umin{s}{T(e, x, s)})
\]
para todo~$x$.
\end{thm}

\begin{proof}[Esboço da demonstração]
Para qualquer modelo de computação, pode-se definir rigorosamente uma
descrição da função computável~$f$ e codificar tal descrição por meio
de um número natural~$e$. Também se pode definir rigorosamente uma noção
de ``sequência de computação'' que registre o processo de calcular a
função de índice~$e$ na entrada~$x$. Essa sequência de computação pode
igualmente ser codificada por um número~$s$. Isso pode ser feito de tal
modo que
\begin{enumerate}
  \item a relação $T(e, x, s)$, que vale se, e somente se, um número~$s$
  codifica a sequência de computação da função de índice~$e$ na
  entrada~$x$, e
  \item a função $U(s)$, que associa uma sequência de computação
  codificada por~$s$ ao resultado final dessa computação,
\end{enumerate}
sejam ambas computáveis. Na verdade, a relação~$T$ e a função~$U$
são recursivas primitivas.
\end{proof}

\begin{explain}
Para dar uma demonstração rigorosa do Teorema da Forma Normal, teríamos
de fixar um modelo de computação, realizar em detalhe a codificação das
descrições das funções computáveis e das sequências de computação, e
verificar que a relação~$T$ e a função~$U$ são recursivas primitivas.
Para a maioria das aplicações, basta que $T$ e~$U$ sejam computáveis
e que $U$~seja total.

Talvez a melhor maneira de lembrar a demonstração do teorema da forma
normal seja como um lema: $\umin{s}{T(e, x, s)}$ procura uma sequência
de computação da função de índice~$e$ na entrada~$x$, e $U$ devolve a
saída da sequência de computação, caso seja possível encontrar uma.
\end{explain}

Se o modelo de computação for o das funções recursivas parciais (que é
o que Kleene usou originalmente), o teorema mostra que apenas um único
uso da busca ilimitada, isto é, um único operador
$\umin{y}{f(\vec x, y)}$, é necessário para definir qualquer função.
Nesse sentido, ele mostra que toda função recursiva parcial tem uma
forma normal.

$T$ e $U$ podem ser usados para definir a enumeração $\cfind{0}$,
$\cfind{1}$, $\cfind{2}$, \dots. De agora em diante, suporemos que
fixamos uma escolha adequada de $T$ e~$U$ e tomaremos a equação
\[
\cfind{e}(x) \simeq U(\umin{s}{T(e,x,s)})
\]
como a \emph{definição} de $\cfind{e}$.

Eis outro fato útil:

\begin{thm}
Toda função parcialmente computável tem infinitos índices.
\end{thm}

Isso também é intuitivamente claro. Dada qualquer (descrição de uma)
função computável, pode-se produzir uma descrição diferente que calcule
a mesma função (o mesmo par de entrada e saída), mas o faça, por
exemplo, realizando primeiro algo que não tenha efeito sobre a
computação (digamos, testar se $0 = 0$, contar até $5$ etc.). O índice
da descrição alterada será sempre diferente do índice original. Ambos
são índices da mesma função, apenas calculada de maneiras ligeiramente
diferentes.

\end{document}

